Language-model agents are deployed for narrow jobs that are specified,
in practice, by a few dozen example queries, a task-level success
check, and metered access to a stronger teacher that charges for every
token. We study few-shot agent specialization under exactly these
terms. Our first finding is negative. Across seven current
distillation baselines spanning the labeling, feedback,
self-knowledge, and on-policy branches, every method that uses teacher
tokens as training targets ends below the untrained student.
Maximum-likelihood imitation provably descends the divergence between
teacher and student behavior distributions, and we find empirically
that the induced task-extraneous drift offsets useful transfer while
eroding abilities the student already has. We therefore propose
budgeted verifier-guided self-distillation. Purchased demonstrations
enter only the student's context, where they raise the probability
that the student discovers successful behavior by itself, by an order
of magnitude in our experiments, and the verifier determines how the
student's own trajectories contribute to learning. A gradient-space
representation of task demand prices each candidate purchase by the
self-supervision it is expected to unlock per teacher token, and
spending stops endogenously when no candidate retains positive
utility. Under matched budgets on agent benchmarks, our method is the
only one that consistently improves on its base student.
\todo{finalize headline numbers at freeze.}
